\section{Computational Protocol}\label{app:computational-protocol}

\subsection{Reference environment}

The public proof artifact is the immutable Git snapshot

\begin{center}
\url{https://github.com/whzy3185/math/tree/c81be34a3b12a7ac47adbb4499c475df7bf4fc04}
\end{center}

and its submission binding is
\path{research/reproducibility/}\allowbreak\path{target_a_submission_}\allowbreak\path{artifact_manifest.json}. The
manifest records the artifact roles and SHA-256 hashes for Theorems~\ref{thm:smallest-counterexample} and~\ref{thm:low-period-frontier}.
Its verifier reads the named files from the pinned Git object, rather than
silently accepting mutable working-tree copies. The snapshot contains all
checker source, the complete 2,626-row low-period certificate table, the
order-24 through order-30 checkpoint data, compact minimality artifacts,
regeneration summary, and pinned dependencies. A DOI-bearing archival release
is desirable at submission; the full commit identifier already makes this
draft's public artifact immutable and content-addressed.

The reference interpreter and packages are:

\begin{center}
\begin{tabular}{ll}
\toprule
Component & Version \\
\midrule
Python & 3.12.13 \\
NumPy & 2.3.5 \\
SymPy & 1.14.0 \\
pytest & 9.1.1 \\
\bottomrule
\end{tabular}
\end{center}

From a clean checkout with Python 3.12 available, a portable environment is
created by

\begin{lstlisting}[language=bash]
python3.12 -m venv .venv-target-a
.venv-target-a/bin/python -m pip install --upgrade pip
.venv-target-a/bin/python -m pip install \
  -r research/reproducibility/requirements-target-a.txt
\end{lstlisting}

All commands below are repository-relative.

\subsection{Fast verification}

The complete default regression suite is run by

\begin{lstlisting}[language=bash]
.venv-target-a/bin/python -m pytest -q research/scripts
\end{lstlisting}

The principal theorem checkers may also be run separately:

\begin{lstlisting}[language=bash]
.venv-target-a/bin/python research/scripts/verify_target_a_minimality_certificate.py
.venv-target-a/bin/python research/scripts/verify_target_a_period8_infinite_family.py
.venv-target-a/bin/python research/scripts/verify_target_a_period8_sharp_constant.py
.venv-target-a/bin/python research/scripts/verify_target_a_period8_structural_mechanism.py
.venv-target-a/bin/python research/scripts/verify_target_a_general_period_moments.py
.venv-target-a/bin/python research/scripts/verify_target_a_low_period_spectral_frontier.py
.venv-target-a/bin/python research/scripts/verify_target_a_low_period_structural_frontier.py
.venv-target-a/bin/python research/scripts/verify_target_a_periodic_operator_equivalences.py
.venv-target-a/bin/python research/scripts/verify_target_a_submission_artifact_manifest.py
\end{lstlisting}

Each command fails closed on a mismatched dependency, count, exact
certificate, scope flag, or source digest.

\subsection{Slow generator audits}

The strongest representative-set audit is run directly by

\begin{lstlisting}[language=bash]
.venv-target-a/bin/python \
  research/scripts/target_a_record_set_audit.py \
  --n 24 26 28 30
\end{lstlisting}

The driver builds the independent C scanner, streams the production FKM
representatives into a disk-mapped table, and requires exact destructive set
consumption. It reports zero missing, duplicate, remaining, canonicality, and
orbit-size errors, together with exact defect and orbit-size histograms. This
audit reconstructs the complete representative set; it does not recompute the
spectral decision attached to every state.

The independent large-order spectral decisions are rerun by

\begin{lstlisting}[language=bash]
.venv-target-a/bin/python \
  research/scripts/target_a_independent_spectral_audit.py \
  --n 24 26 28 30
\end{lstlisting}

This command consumes records emitted by the C full-space scan, independently
reconstructs both holonomy matrices, and verifies every nonoptimizer with an
exact integer Rayleigh quotient. It is a full decision rerun rather than a
checkpoint replay. The recorded reference execution took tens of minutes; its
per-order counts, independent certificate digests, source hashes, and zero
uncertified-state checks are bound by
\path{research/reproducibility/}\allowbreak
\path{target_a_computational_evidence_manifest.json}.

The three generator audits skipped by the default suite are enabled explicitly:

\begin{lstlisting}[language=bash]
TARGET_A_RUN_SLOW_GENERATOR_TESTS=1 \
.venv-target-a/bin/python -m pytest -q \
  research/scripts/test_target_a_direct_bracelets.py::DirectBraceletTests::test_direct_generator_burnside_n26 \
  research/scripts/test_target_a_direct_bracelets.py::DirectBraceletTests::test_direct_generator_burnside_n28 \
  research/scripts/test_target_a_direct_bracelets.py::DirectBraceletTests::test_direct_generator_burnside_n30
\end{lstlisting}

These tests independently check shell totals, Burnside counts, represented
class totals, ordering, and parity. They are audits, not fresh spectral
regeneration.

\subsection{Regeneration versus integrity replay}

Fresh regeneration starts from the canonical generator, reconstructs every
spectral state, recomputes its exact decision, and writes a new checkpoint
chain. Integrity replay reads an existing chain and verifies its hashes,
counts, cursor, and aggregate records. Both are useful, but only the former
re-executes the mathematical decision for every state.

For the production orders, choose an empty external directory and run

\begin{lstlisting}[language=bash]
export OUT="${TARGET_A_OUTPUT:?set TARGET_A_OUTPUT}"
mkdir -p "$OUT/checkpoints" "$OUT/results"
for n in 24 26 28 30; do
  .venv-target-a/bin/python research/scripts/target_a_minimality_search.py \
    --n "$n" \
    --checkpoint-dir "$OUT/checkpoints/n$n" \
    --output "$OUT/results/target_a_search_n$n.json"
done
\end{lstlisting}

An interrupted order is resumed by adding \(--resume\) to its command. Each
result JSON records the completed state and shell totals, ordered input and
certificate hashes, terminal checkpoint-chain hash, elapsed time, and peak
resident memory. A successful reference run on an Apple M5 using Python
3.12.13, NumPy 2.3.5, and SymPy 1.14.0 produced:

\begin{table}[tbp]
\centering
\caption{Recorded regeneration resources and terminal checkpoint hashes.}
\label{tab:regeneration-resources}
\begingroup
\renewcommand{\arraystretch}{1.12}
\scriptsize
\resizebox{\textwidth}{!}{%
\begin{tabular}{lllllll}
\toprule
\(n\) & states & chunks & time (s) & peak RSS (MiB) & checkpoint disk & terminal chain SHA-256 \\
\midrule
24 & 353,812 & 28 & 25.3 & 75.1 & 1.1 MiB & \texttt{\detokenize{2dde869aea5da4f040e67a4fef3e93b5f35f5fb42d75f6d48820439f418c83c1}} \\
26 & 1,299,064 & 76 & 117.2 & 84.4 & 3.1 MiB & \texttt{\detokenize{c515350b8bea840c04448086fbc98523615364c05ca837553c93efb933bc0c4e}} \\
28 & 4,810,472 & 250 & 414.9 & 86.7 & 12 MiB & \texttt{\detokenize{7ba200b05590b2a9c0ea1121f25f89ddf1d294d0c043417e69f78f764f6e8ee1}} \\
30 & 17,929,600 & 908 & 1542.3 & 135.2 & 43 MiB & \texttt{\detokenize{b7fd264eece645eead187424152ae810a9ff940e37ffc5649b5ddf65aa31d59d}} \\
\bottomrule
\end{tabular}
}
\endgroup
\end{table}

Runtime is hardware-dependent; the observed sequential total was about 35
minutes, peak per-process memory stayed below 150 MiB, and the four checkpoint
directories occupied about 60 MiB. The compact committed reproduction summary
also records the ordered input and ordered certificate hashes, so a fresh run
can be compared without access to an off-repository timing log. The committed
chains themselves are replayed by

\begin{lstlisting}[language=bash]
.venv-target-a/bin/python research/scripts/target_a_checkpoint_replay.py \
  --n 24 26 28 30 \
  --output "$OUT/target_a_checkpoint_replay.json"
\end{lstlisting}

The production checkpoints are committed. The fresh regeneration chunks and
full runtime logs are retained outside the repository; their compact manifest
records state counts, chunk counts, ordered input and certificate digests, and
terminal chain digests. Agreement of these logical fingerprints, rather than
byte equality of timing metadata, defines mismatch zero.

\subsection{Negative tests}

Regression tests alter one logical field at a time and require rejection.
They cover, among other faults, missing orbit representatives, duplicated
canonical words, nondeterministic orbit identifiers, dependency digest drift,
false primitive periods, altered Rayleigh numerators, the lower radical branch,
incorrect moment direction, numeric previews promoted to proof, and forbidden
all-period scope. This adversarial layer is important because many positive
checks share the same exact arithmetic primitives.
